\begin{center}\begin{tikzpicture}[font=\small,>=stealth,line join=round]
\draw[gray!40,rounded corners=3pt] (0,0) rectangle (14.2,8.3);\node at (7.1,7.700000000000001) {例3：递减函数的两点，高低与左右顺序相反};
\begin{scope}[shift={(1.6454545454545455,2.161111111111111)},x=1.3636363636363635cm,y=1.0222222222222221cm]\draw[->,gray] (-0.4,0)--(4,0) node[right]{$x$};\draw[->,gray] (0,-0.5)--(0,4) node[above]{$y$};\draw[orange!85!black,thick,solid,domain=0.27:3.5,samples=180] plot(\x,{1/\x});\fill[orange!85!black] (1,1) circle(2.2pt);\fill[orange!85!black] (2,0.5) circle(2.2pt);\node[above right=4pt] at (1,1){$(a,f(a))$};\node[above right=4pt] at (2,.5){$(b,f(b))$};\end{scope}\node[anchor=west,align=left,text width=5.0cm] at (8.4,4.0) {图取$\alpha=-1$\\[9pt]$a=1<b=2$\\[9pt]却有$f(a)=1>f(b)=\frac12$\\[9pt]一般地，$\alpha<0$时\\[9pt]$a^\alpha>b^\alpha\Rightarrow a<b$};
\node[text width=12.6cm,align=center] at (7.1,.65) {图给出一组代表点；正文的反证推理适用于任意负指数与正数a、b。};
\end{tikzpicture}\end{center}
